\newpage

\section{Interfaccia Utente}


\subsection{FSM e Contratto view\_interface\_t}

Il sistema UI è una FSM dove ogni stato corrisponde a una schermata.
\texttt{view\_id\_t} enumera le view disponibili; il puntatore
\texttt{current\_view} in \texttt{ui\_manager.cpp} punta all'istanza
\texttt{view\_interface\_t} attiva. Le transizioni avvengono solo tramite
\texttt{ui\_manager\_navigate\_to()} che chiama la funzione \texttt{on\_exit} della
view uscente, aggiorna il puntatore, attiva il feedback sonoro di navigazione
e infine esegue la funzione \texttt{on\_enter} della nuova view.

\begin{lstlisting}[caption={view\_definitions.h --- contratto funzionale di una view}]
typedef struct {
    void (*on_enter)(void);
    void (*on_update)(const void* data);
    void (*on_input)(button_t btn, button_state_t state);
    void (*on_exit)(void);
} view_interface_t;
\end{lstlisting}


\subsection{Rendering e Timing OLED}

Il ciclo di render segue il pattern \texttt{clearDisplay()} $\to$ scrittura
nel frame buffer $\to$ \texttt{display()}. La chiamata \texttt{display()}
trasferisce 1024~byte (128$\times$64~bit) via I2C a 400~kHz: circa 25~ms di
trasferimento. Per questo \texttt{task\_ui} è programmato ogni 33~ms,
evitando di accodare un secondo trasferimento prima che il precedente sia
completato.

\subsection{View Implementate}

\textbf{view\_home} implementa un menu a cursore con navigazione ciclica.
La riga selezionata è renderizzata in modalità inversa
(rettangolo bianco riempito + testo nero), tecnica che non richiede bitmap.
La mappatura tra posizione del cursore e destinazione è una lookup table:
\begin{lstlisting}
view_id_t targets[] = {VIEW_ID_GPS, VIEW_ID_METEO, VIEW_ID_SETTINGS};
ui_manager_navigate_to(targets[menu_cursor]);
\end{lstlisting}

\textbf{view\_gps} converte la velocità (\texttt{speed\_ms * 3.6f}),
formatta le coordinate con 6 decimali e mostra il conteggio satelliti con
padding fisso a 2 cifre (\texttt{\%02d}).

\textbf{view\_meteo} legge il campo \texttt{weather.valid} come sentinella.
La descrizione meteo è derivata con un sistema a soglie da \texttt{weather\_code}:
0~=~Clear~Sky, 1--3~=~Partly~Cloudy, $\geq$60~=~Rainy. Le soglie seguono
la classificazione WMO usata da OpenMeteo.

\textbf{view\_settings} gestisce 5 voci con ciclo della luminosità
secondo la formula:
\begin{lstlisting}
global_cfg.oled_brightness =
    (global_cfg.oled_brightness >= 250)
    ? 50
    : ((global_cfg.oled_brightness / 50) + 1) * 50;
\end{lstlisting}
Il risultato è una sequenza circolare 50 $\to$ 100 $\to$ 150 $\to$ 200
$\to$ 250 $\to$ 50 con incrementi di 50. L'effetto è applicato
immediatamente su hardware tramite
\texttt{display\_set\_brightness(global\_cfg.oled\_brightness)}

\textbf{view\_info} mostra versione firmware e uptime, formattato HH:MM:SS con divisioni intere su \texttt{uptime\_s}.

\subsubsection*{Overlay Errori}

\texttt{ui\_manager\_update()} controlla \texttt{data->flags.error\_active}
prima di delegare alla view normale. Se attivo e non confermato,
\texttt{render\_error\_overlay()} sovrascrive il buffer mostrando categoria
e codice in esadecimale. Al primo input successivo,
\texttt{ui\_manager\_dispatch\_input()} imposta \texttt{error\_acknowledged}
e chiama \texttt{sys\_manager\_clear\_error()}, ripristinando la navigazione.


\subsubsection*{Feedback Acustico}

I beep sono differenziati per tipo di interazione in
\texttt{ui\_manager\_dispatch\_input()}: 2400~Hz/15~ms per navigazione
(UP/DOWN), 1800~Hz/20~ms per BACK. La navigazione verso una nuova view
genera 1500~Hz/80~ms. 
 
